\section{表示的检出能力：形状变化与平移负对照}\label{sec:h1}

表~\ref{tab:h1}已完整报告两个场景的尺寸校正功效与场景差。形状变化下 BPI 的功效接近 $1$、MOMENTS 约 $0.06$，差异接近一个单位；平移负对照则正好反转，BPI 接近名义水平而 MOMENTS 接近 $1$，两个场景差的区间均不跨零。圆到八字形构造改变潜在几何并匹配单时点一、二阶矩；检出差异可能同时涉及连通、循环和度量几何变化，本实验未分离一维同调的贡献。

\begin{table}[H]
\centering
\small
\caption{两类备择下四个预设条件的等权尺寸校正功效与场景差（BPI 减 MOMENTS）。}\label{tab:h1}
\begin{tabular}{lccc}
\toprule
场景 & BPI 功效 & MOMENTS 功效 & 差值 $[95\%$ CI$]$\\
\midrule
形状变化 & 1.00000 & 0.06203 & $0.93797\ [0.92969,0.94609]$\\
均值平移 & 0.05375 & 1.00000 & $-0.94625\ [-0.95344,-0.93844]$\\
\bottomrule
\end{tabular}

\vspace{0.2em}
\begin{minipage}{0.95\textwidth}
\footnotesize 注：功效以各自匹配的零假设 $0.95$ 经验临界值计算；差值为两表示场景平均功效之差，其 $95\%$ 区间由 2000 次配对双层整簇自助法获得。逐条件的 nMAE、联合成功率与名义功效见补充材料 S7。
\end{minipage}
\end{table}


平移负对照有直接的结构解释。非重叠条件 $s=W=64$、$M=192$ 时 $T=12288$、$\tau=6144=96W$，变点恰落在第 96 与第 97 个窗口之间，没有窗口同时包含变化前后的观测；将 $H_1$ 与 $H_0$ 耦合到相同的潜变量与噪声路径时，每个后段窗口只是整体平移，故 $Y_{i,H_1}^{\mathrm{BPI}}=Y_{i,H_0}^{\mathrm{BPI}}$ 且 $T_{M,H_1}^{\mathrm{BPI}}=T_{M,H_0}^{\mathrm{BPI}}$。在该非重叠且边界对齐的构造中，BPI 统计量与匹配零假设具有相同分布，其检出概率等于匹配零假设下该经验临界值对应的拒绝概率；有限样本估计接近 $5\%$，但不要求精确等于 $0.05$，该结论对两种 $\rho$ 均成立。重叠条件 $s=16$、$M=192$ 时 $T=3120$、$\tau=1560$，只有第 95--98 号窗口跨越变点，其余窗口的 BPI 不变，故差异只能来自这些混合窗口。由于欧氏距离 Vietoris--Rips 构造对整体平移不变，该场景只能作为表示不变性的负对照，不能用于评价一般平移备择下的检出性能。

四个条件各 $1600$ 次重复；场景平均功效为 $1.00000$ 意味着四格估计均为 $1$，不构成总体功效恒等于 $1$ 的概率保证。定位与联合成功率见补充材料 S7：场景汇总的 nMAE 为 $0.00330$--$0.20255$，联合成功率为 $0.0009$--$0.9839$。本文不将其作为主贡献：位置误差由观测扫描曲线决定，校准方法只通过拒绝事件影响联合成功率，而失校准条件下的较高联合成功率可能部分来自过高的拒绝概率，须与匹配的零假设尺寸共同解释。
